\documentclass[11pt]{article}
\usepackage[a4paper,margin=28mm]{geometry}
\usepackage{amsmath,amssymb,amsthm,mathtools}
\usepackage[hidelinks]{hyperref}

\newtheorem{theorem}{Theorem}
\newtheorem{proposition}{Proposition}
\newtheorem{lemma}{Lemma}
\newtheorem{corollary}{Corollary}
\theoremstyle{remark}
\newtheorem{remark}{Remark}

\newcommand{\Hh}[2]{H_{#1}^{(#2)}}
\newcommand{\ZZ}{\mathbb Z}
\newcommand{\QQ}{\mathbb Q}
\DeclareMathOperator{\Res}{Res}

\title{The compact Brown--Zudilin weights:\\ proofs for the companion rows}
\author{}
\date{July 27, 2026}

\begin{document}
\maketitle

\begin{abstract}
Let $T(n,k,l)=\binom{n+k}{n}\binom nk^2\binom{n+l}{n}\binom nl^2\binom{n+k+l}{n}$
and let $Q_n,\hat P_n,P_n$ be the coefficient sequences of the totally symmetric
Brown--Zudilin cellular linear forms for $\zeta(5)$, so that
$I_n=Q_n(2\zeta(5)+4\zeta(3)\zeta(2))-4\hat P_n\zeta(2)-2P_n$.
We prove the compact closed form
$\hat P_n=\sum_{k,l}T(n,k,l)\,\hat w_3(n,k,l)$ for the middle companion row,
and we prove the exact Barnes-rational representation
$P_n=-\tfrac12\sum_{k,l}T(n,k,l)\,W^{(1)}(k,l)$
for the top row, where $W^{(1)}$ is an explicit polynomial in harmonic and
finite Euler sums.  Both proofs are uniform in $n$ and use no
$\QQ$-linear-independence assumptions about zeta values.  The remaining step
from the Barnes-rational representation to the compact three-term weight
$w_5$ is a single identity between finite sums, which we isolate, reduce to a
linear-algebra problem over a space of proved residue functionals, and
discuss in the final section.
\end{abstract}

\section{Statement of results}\label{sec:results}

Throughout, $\Hh mr=\sum_{i=1}^m i^{-r}$ (with $\Hh mr=0$ for $m\le 0$), and
\[
 T(n,k,l)
 =\binom{n+k}{n}\binom nk^2\binom{n+l}{n}\binom nl^2\binom{n+k+l}{n},
 \qquad 0\le k,l\le n .
\]
Brown--Zudilin \cite[Sec.~2]{BZ} attach to the totally symmetric cellular
integral
\[
 I_n=\idotsint\limits_{0<t_1<\dots<t_5<1}
 \Bigl(\tfrac{t_1(t_2-t_1)(t_3-t_2)(t_4-t_3)(t_5-t_4)(1-t_5)}
 {(t_3-t_1)t_3(1-t_4)(t_4-t_2)(t_5-t_2)}\Bigr)^{\!n}
 \frac{dt_1\cdots dt_5}{(t_3-t_1)t_3(1-t_4)(t_4-t_2)(t_5-t_2)}
\]
the decomposition
\begin{equation}\label{eq:In}
 I_n=Q_n\bigl(2\zeta(5)+4\zeta(3)\zeta(2)\bigr)-4\hat P_n\,\zeta(2)-2P_n,
\end{equation}
where $Q_n,\hat P_n,P_n$ are rational numbers.  All four sequences satisfy
one third-order Ap\'ery-type recursion (produced by creative telescoping for
$I_n$ and checked directly for the rational sequences), with initial data
\[
 Q_0=1,\ Q_1=21,\ Q_2=2989,\qquad
 \hat P_0=0,\ \hat P_1=\tfrac{101}4,\ \hat P_2=\tfrac{344923}{96},\qquad
 P_0=0,\ P_1=\tfrac{87}4,\ P_2=\tfrac{1190161}{384},
\]
and \eqref{eq:In} is \emph{proved} in \cite{BZ} by exactly this route:
both sides satisfy the same recursion and agree at $n=0,1,2$.  In
particular \eqref{eq:In} is an identity between real numbers that does not
presuppose any linear independence of zeta values; the sequences
$Q_n,\hat P_n,P_n$ are pinned by the recursion and the initial data.
Brown--Zudilin also prove
\begin{equation}\label{eq:Qn}
 Q_n=\sum_{k,l=0}^{n}T(n,k,l),
\end{equation}
again by checking the recursion.

Define, for $r\ge 1$,
\[
 A_r(x)=\Hh{n+x}r-\Hh xr,\qquad
 B_r(x)=\Hh{n-x}r-\Hh xr,\qquad
 C_r=\Hh{n+k+l}r-\Hh{k+l}r,
\]
\[
 \alpha=A_1(k)-A_1(l),\qquad \beta=B_1(k)-B_1(l),\qquad \Psi=\tfrac\alpha2+\beta,
\]
and the \emph{compact weights}
\[
 \hat w_3=\Hh{n+k}3-\Psi\,\Hh{n+k}2,
\qquad
 w_5=\Hh{n+k}5+\tfrac{\alpha-\beta}2\Hh{n+k}4
 +\Bigl(\tfrac{A_2(k)+A_2(l)}4-\tfrac{\alpha\Psi}2\Bigr)\Hh{n+k}3 .
\]
Finally, let $L_k=-A_1(k)-C_1-2B_1(k)$ and $L_l=-A_1(l)-C_1-2B_1(l)$, and let
$r_{11},r_{12},r_{21},r_{22}$ be the explicit weight-$\le5$ expressions in
harmonic and finite Euler sums displayed in Proposition~\ref{prop:ratparts}
below.  Put
\begin{equation}\label{eq:W1}
 W^{(1)}(k,l)=r_{22}(k,l)+L_k\,r_{12}(k,l)+L_l\,r_{21}(k,l)
 +(L_kL_l-C_2)\,r_{11}(k,l).
\end{equation}

The results of this paper are:

\begin{theorem}[middle row]\label{thm:middle}
For every $n\ge0$,
\[
 \hat P_n=\sum_{k,l=0}^{n}T(n,k,l)\,\hat w_3(n,k,l).
\]
\end{theorem}

\begin{theorem}[top row, Barnes-rational form]\label{thm:anchor}
For every $n\ge0$,
\[
 P_n=-\frac12\sum_{k,l=0}^{n}T(n,k,l)\,W^{(1)}(k,l).
\]
\end{theorem}

\begin{theorem}[top row, compact form; conditional]\label{thm:top}
If the bridge identity
\begin{equation}\label{eq:T3}
 \sum_{k,l=0}^{n}T(n,k,l)\,W^{(1)}(k,l)
 =-2\sum_{k,l=0}^{n}T(n,k,l)\,w_5^{\mathrm{sym}}(n,k,l)
 \qquad(n\ge0)
\end{equation}
holds, where $w_5^{\mathrm{sym}}$ is the $k\leftrightarrow l$ symmetrisation
of $w_5$, then
\[
 P_n=\sum_{k,l=0}^{n}T(n,k,l)\,w_5(n,k,l)\qquad (n\ge 0).
\]
Identity \eqref{eq:T3} has been verified exactly for $n\le12$ and modulo two
$31$-bit primes for $n\le80$; Section~\ref{sec:T3} reduces it to membership
of one explicit vector in the span of proved residue functionals and records
the current state of that reduction.  It is the only unproved statement in
this paper.
\end{theorem}

Everything else below is a complete proof of
Theorems~\ref{thm:middle} and~\ref{thm:anchor}.  The logical skeleton is:
\begin{enumerate}
\item the Barnes kernel and its bivariate partial fractions
      (Sections~\ref{sec:kernel}--\ref{sec:pf});
\item the recombination $I_n=\sum T\cdot W_B$ over four universal
      sine-kernel integrals (Section~\ref{sec:recomb});
\item the exact evaluation of the universal integrals in the basis
      $1,\zeta(2),\zeta(3),\zeta(4),\zeta(5),\zeta(2)\zeta(3)$
      (Section~\ref{sec:universal});
\item four kernel identities: the $\zeta(5)$ and $\zeta(2)\zeta(3)$
      coefficients equal $2Q_n$ and $4Q_n$; the $\zeta(4)$ and $\zeta(3)$
      coefficients vanish; the $\zeta(2)$ coefficient equals
      $-4\sum T\,w_3^{\mathrm{sym}}$ (Section~\ref{sec:kernelids});
\item the middle row by descent through the generalised Beukers integral
      and Zudilin's evaluation lemma (Section~\ref{sec:middle});
\item the subtraction anchor: combining 1--5 with \eqref{eq:In} kills every
      transcendental term exactly and leaves Theorem~\ref{thm:anchor}
      (Section~\ref{sec:anchor}).
\end{enumerate}

\begin{remark}
Step 6 is where previous attempts stalled, because comparing coefficients of
$\zeta(2)$ or $\zeta(5)$ between two evaluations of $I_n$ is illegitimate
while the $\QQ$-linear independence of
$1,\zeta(2),\zeta(3),\zeta(5),\zeta(2)\zeta(3)$ remains open.  The
subtraction route avoids this: once the middle row is proved
\emph{independently} (Theorem~\ref{thm:middle}), the two evaluations of the
real number $I_n$ differ only in their rational terms, and equality of real
numbers then forces equality of those rational terms.  No independence is
used anywhere.
\end{remark}

\section{The Barnes kernel}\label{sec:kernel}

Brown--Zudilin's Barnes-type representation \cite[eq.~(intJ)]{BZ},
specialised to the totally symmetric parameters
$p_0=p_1=p_2=p_4=p_5=p_6=q_1=\dots=q_5=n$, $p_3=2n$ (for which the
normalising factorial prefactor collapses to $n!$), reads
\begin{equation}\label{eq:barnes}
 I_n=\frac{n!}{(2\pi i)^2}
 \int_{-c_1-i\infty}^{-c_1+i\infty}\!\!\int_{-c_2-i\infty}^{-c_2+i\infty}
 \frac{\Gamma(n{+}1{+}s)^3\,\Gamma(-s)}{\Gamma(2n{+}2{+}s)^2}\,
 \frac{\Gamma(n{+}1{+}t)^3\,\Gamma(-t)}{\Gamma(2n{+}2{+}t)^2}\,
 \Gamma(2n{+}2{+}s{+}t)\,\Gamma(-n{-}1{-}s{-}t)\,ds\,dt,
\end{equation}
valid for $0<c_1,c_2<n+1$ and $c_1+c_2>n+1$; we fix
$c_1=c_2=n+\tfrac23$.

\begin{lemma}\label{lem:kernel}
With $c_1=c_2=n+\frac23$,
\[
 I_n=\frac{(-1)^n\,n!}{(2\pi i)^2}
 \iint R_n(s,t)\,
 \frac{\pi}{\sin\pi s}\,\frac{\pi}{\sin\pi t}\,\frac{\pi}{\sin\pi(s+t)}
 \,ds\,dt,
\]
where
\[
 R_n(s,t)=
 \frac{\prod_{j=1}^n(s+j)\,\prod_{j=1}^n(t+j)\,\prod_{j=n+2}^{2n+1}(s+t+j)}
      {\prod_{j=n+1}^{2n+1}(s+j)^2\,\prod_{j=n+1}^{2n+1}(t+j)^2}.
\]
\end{lemma}

\begin{proof}
Apply the reflection formula $\Gamma(z)\Gamma(1-z)=\pi/\sin\pi z$ three
times:
\[
 \Gamma(-s)=\frac{-\pi}{\sin(\pi s)\,\Gamma(1+s)},\qquad
 \Gamma(-t)=\frac{-\pi}{\sin(\pi t)\,\Gamma(1+t)},
\]
and, using $\sin\pi(-n-1-u)=(-1)^{n}\sin\pi u$,
\[
 \Gamma(-n-1-s-t)=\frac{(-1)^n\,\pi}{\sin(\pi(s+t))\,\Gamma(n+2+s+t)} .
\]
The three sign factors multiply to $(-1)^n$.  The remaining gamma quotients
are rational:
\[
 \frac{\Gamma(n+1+s)^3}{\Gamma(2n+2+s)^2\,\Gamma(1+s)}
 =\frac{\prod_{j=1}^n(s+j)}{\prod_{j=n+1}^{2n+1}(s+j)^2},
 \qquad
 \frac{\Gamma(2n+2+s+t)}{\Gamma(n+2+s+t)}
 =\prod_{j=n+2}^{2n+1}(s+t+j). \qedhere
\]
\end{proof}

On the chosen contours the integrand of Lemma~\ref{lem:kernel} is absolutely
integrable: on vertical lines with non-integral real part each factor
$\pi/\sin$ decays like $e^{-\pi|\Im|}$ in its argument, so the triple decays
like $e^{-\pi(|y_1|+|y_2|+|y_1+y_2|)}\le e^{-2\pi\max(|y_1|,|y_2|)}$, while
$R_n$ has polynomial decay $O(|s|^{-2}|t|^{-2})$ in each variable.  All the
rearrangements below happen inside this absolutely convergent double
integral.

\section{Bivariate partial fractions of the kernel}\label{sec:pf}

\begin{proposition}\label{prop:pf}
The rational function $R_n$ has poles exactly on the grid
$s=s_k:=-n-1-k$, $t=t_l:=-n-1-l$, $0\le k,l\le n$, of order at most two in
each variable, and
\[
 R_n(s,t)=\sum_{k,l=0}^{n}\left[
 \frac{C_{22}}{(s-s_k)^2(t-t_l)^2}
 +\frac{C_{12}}{(s-s_k)(t-t_l)^2}
 +\frac{C_{21}}{(s-s_k)^2(t-t_l)}
 +\frac{C_{11}}{(s-s_k)(t-t_l)}\right],
\]
with local data (depending on $n,k,l$)
\[
 (-1)^n n!\,C_{22}=T(n,k,l),\qquad
 \frac{C_{12}}{C_{22}}=L_k,\qquad
 \frac{C_{21}}{C_{22}}=L_l,\qquad
 \frac{C_{11}}{C_{22}}=L_kL_l-C_2 .
\]
\end{proposition}

\begin{proof}
For fixed generic $t$, $R_n(\cdot,t)$ is a proper rational function
($O(s^{-2})$ at infinity) with double poles at $s_0,\dots,s_n$ only, so it
equals the sum of its principal parts; expanding each principal-part
coefficient (a rational function of $t$ of the same shape) in $t$ gives the
displayed grid decomposition with no polynomial remainder.

For the local data, write $s=s_k+x$, $t=t_l+y$ and factor the germ as
$R_n=\Lambda\,F(x,y)/(x^2y^2)$ with $F(0,0)=1$.  Then
$C_{22}=\Lambda$, $C_{12}/C_{22}=\partial_xF(0,0)$,
$C_{21}/C_{22}=\partial_yF(0,0)$,
$C_{11}/C_{22}=\partial_x\partial_yF(0,0)$.

\emph{Leading coefficient.}
$\lim_{s\to s_k}(s-s_k)^2\prod_{j=1}^n(s+j)\big/\prod_{j=n+1}^{2n+1}(s+j)^2$:
the numerator at $s_k$ is $\prod_{j=1}^n(j-n-1-k)=(-1)^n(n+k)!/k!$, while
the surviving denominator factors are
$\prod_{j\ne n+k+1}(j-n-1-k)=(-1)^kk!\,(n-k)!$, squared.  Hence the $s$-factor
contributes $(-1)^n(n+k)!/\bigl(k!^3((n-k)!)^2\bigr)$, and likewise in $t$.
Finally $\prod_{j=n+2}^{2n+1}(s_k+t_l+j)=(-1)^n(n+k+l)!/(k+l)!$.  Multiplying,
\[
 \Lambda=(-1)^n\,
 \frac{(n+k)!\,(n+l)!\,(n+k+l)!}{k!^3\,l!^3\,(k+l)!\,((n-k)!)^2((n-l)!)^2}
 =\frac{(-1)^n}{n!}\,T(n,k,l).
\]

\emph{First jets.}  $\partial_x\log F(0,0)$ collects the logarithmic
derivatives at $s=s_k$ of the three $s$-dependent products:
$\sum_{j=1}^n\frac1{s_k+j}=-\sum_{i=k+1}^{k+n}\frac1i=-A_1(k)$;
$-2\sum_{j\ne n+k+1}\frac1{s_k+j}=-2(-H_k+H_{n-k})=-2B_1(k)$;
and $\sum_{j=n+2}^{2n+1}\frac1{s_k+t_l+j}=-\sum_{i=k+l+1}^{k+l+n}\frac1i=-C_1$.
Hence $\partial_xF(0,0)=-A_1(k)-2B_1(k)-C_1=L_k$, and symmetrically
$\partial_yF(0,0)=L_l$.

\emph{Mixed jet.}  Since only the coupled product depends on $s+t$,
$\partial_x\partial_y\log F(0,0)
=-\sum_{j=n+2}^{2n+1}(s_k+t_l+j)^{-2}=-C_2$, so
$\partial_x\partial_yF(0,0)=L_kL_l-C_2$.
\end{proof}

\section{Recombination over the universal integrals}\label{sec:recomb}

For integers $k_1,k_2\ge0$ and $s_1,s_2\in\{1,2\}$ define, following
\cite[Remark~1]{BZ},
\begin{equation}\label{eq:universal}
 I^{(s_1,s_2)}_{k_1,k_2}
 =\frac1{(2\pi i)^2}
 \int_{\frac13-i\infty}^{\frac13+i\infty}\!\!\int_{\frac13-i\infty}^{\frac13+i\infty}
 \frac{1}{(u+k_1)^{s_1}(v+k_2)^{s_2}}\,
 \frac{\pi}{\sin\pi u}\,\frac{\pi}{\sin\pi v}\,\frac{\pi}{\sin\pi(u+v)}
 \,du\,dv .
\end{equation}

\begin{proposition}[recombination]\label{prop:recomb}
For every $n\ge0$,
\[
 I_n=\sum_{k,l=0}^{n}T(n,k,l)\,W_B(n,k,l),
 \qquad
 W_B:=I^{(2,2)}_{k,l}+L_k\,I^{(1,2)}_{k,l}+L_l\,I^{(2,1)}_{k,l}
 +(L_kL_l-C_2)\,I^{(1,1)}_{k,l}.
\]
\end{proposition}

\begin{proof}
Insert the partial fractions of Proposition~\ref{prop:pf} into
Lemma~\ref{lem:kernel}; absolute convergence permits integrating term by
term.  In the $(k,l)$ term substitute $u=s+n+1$, $v=t+n+1$.  The contours
$\Re s=-n-\frac23$, $\Re t=-n-\frac23$ become $\Re u=\Re v=\frac13$, no pole
is crossed, and
\[
 \sin\pi s=(-1)^{n+1}\sin\pi u,\qquad
 \sin\pi t=(-1)^{n+1}\sin\pi v,\qquad
 \sin\pi(s+t)=\sin\pi(u+v),
\]
so the sine triple picks up the factor $(-1)^{2(n+1)}=1$, while
$(s-s_k,t-t_l)=(u+k,v+l)$.  Thus the $(k,l)$ term contributes
$C_{22}I^{(2,2)}_{k,l}+C_{12}I^{(1,2)}_{k,l}+C_{21}I^{(2,1)}_{k,l}
+C_{11}I^{(1,1)}_{k,l}$, and by Proposition~\ref{prop:pf} the prefactor
$(-1)^n n!\,C_{22}=T$ absorbs the $(-1)^n n!$ of Lemma~\ref{lem:kernel}.
\end{proof}

\section{Exact evaluation of the universal integrals}\label{sec:universal}

Brown--Zudilin \cite[Remark~1]{BZ} compute the Fourier--Mellin kernel
\[
 f(u,v)=\frac1{(2\pi i)^2}\iint u^{t_1}v^{t_2}
 \frac{\pi}{\sin\pi t_1}\frac{\pi}{\sin\pi t_2}\frac{\pi}{\sin\pi(t_1+t_2)}
 dt_1\,dt_2
 =
 \begin{cases}
 \dfrac{uv\log u}{(1-u)(1-v)}-\dfrac{u\log(u/v)}{(1-u/v)(1-v)}, & 0<u<v<1,\\[6pt]
 \text{(same with $u\leftrightarrow v$)}, & 0<v<u<1,
 \end{cases}
\]
by residue analysis.  Combining with the Laplace representation
$\frac1{(t+a)^s}=\frac1{\Gamma(s)}\int_0^1 u^{t+a-1}(-\log u)^{s-1}\,du$
(valid on $\Re t=\frac13$, $a\ge0$) gives
\begin{equation}\label{eq:fuv}
 I^{(s_1,s_2)}_{k_1,k_2}
 =\frac{(-1)^{s_1+s_2}}{\Gamma(s_1)\Gamma(s_2)}
 \iint_{[0,1]^2}u^{k_1}v^{k_2}\,f(u,v)\,
 (\log u)^{s_1-1}(\log v)^{s_2-1}\,\frac{du}u\frac{dv}v .
\end{equation}
(The display in \cite{BZ} omits the sign $(-1)^{s_1+s_2}$, i.e.\ writes
$\log$ for $-\log$; the sign is invisible for $(s_1,s_2)=(2,2)$ and $(1,1)$
but reverses $(1,2)$ and $(2,1)$.  It is fixed by the defining contour
integral \eqref{eq:universal}, and independently by direct numerical contour
quadrature and by the calibration
$I^{(2,2)}_{0,0}=I_0=2\zeta(5)+4\zeta(2)\zeta(3)$.)

Expanding \eqref{eq:fuv} on the two triangles by geometric series and
integrating termwise produces an explicit evaluation.  We record it in the
form needed later.  Define, for $x\ge0$ and $r,m\ge1$, the finite Euler sums
\[
 S_{r,m}(x)=\sum_{t=1}^{x}\frac{\Hh tm}{t^{r}},
 \qquad
 U_{r,m}(a,b)=\sum_{t=1}^{a}\frac{\Hh{t+b}m}{t^{r}} .
\]

\begin{proposition}\label{prop:ratparts}
Each $I^{(s_1,s_2)}_{k,l}$ is a $\QQ$-linear combination of
$1,\zeta(2),\zeta(3),\zeta(4),\zeta(5),\zeta(2)\zeta(3)$ whose coefficients
are explicit polynomials in harmonic numbers at $k,l,k+l$ and the finite
Euler sums above.  The coefficients of weight $>s_1+s_2+1$ vanish;
in particular $\zeta(5)$ and $\zeta(2)\zeta(3)$ occur only in $I^{(2,2)}$,
with the constant coefficients
\[
 [\zeta(5)]\,I^{(2,2)}=2,\qquad [\zeta(2)\zeta(3)]\,I^{(2,2)}=4 .
\]
The remaining coefficients are:
\begin{align*}
 [\zeta(4)]:\quad & I^{(1,2)}=I^{(2,1)}=\tfrac{17}4,\qquad
                    I^{(1,1)}=I^{(2,2)}=0;\\[2pt]
 [\zeta(3)]:\quad & I^{(1,1)}=2,\qquad
   I^{(1,2)}=2(H_k-H_{k+l}),\qquad I^{(2,1)}=2(H_l-H_{k+l}),\\
                  & I^{(2,2)}=2\bigl(\Hh k2+\Hh l2-2\Hh{k+l}2\bigr);\\[2pt]
 [\zeta(2)]:\quad & I^{(1,1)}=H_k+H_l-2H_{k+l},\qquad
   I^{(1,2)}=\Hh l2-2\Hh{k+l}2,\qquad I^{(2,1)}=\Hh k2-2\Hh{k+l}2,\\
                  & I^{(2,2)}=-4\Hh{k+l}3;
\end{align*}
and the rational parts $r_{s_1s_2}:=[1]\,I^{(s_1,s_2)}_{k,l}$ are
\begin{align*}
 r_{11}&=(H_{k+l}-H_k-H_l)(\Hh k2+\Hh l2)-\Hh k3-\Hh l3
        +U_{1,2}(k,l)+U_{1,2}(l,k),\\[2pt]
 r_{12}&=-2(H_k+H_l-H_{k+l})\Hh l3
        +\Hh k2\Hh{k+l}2-\tfrac12(\Hh l2)^2+\Hh{k+l}2\Hh l2
        -\tfrac52\Hh l4+2S_{1,3}(l)-U_{2,2}(k,l),\\[2pt]
 r_{21}&=r_{12}\ \text{with}\ k\leftrightarrow l,\\[2pt]
 r_{22}&=-2(\Hh k2+\Hh l2)(\Hh k3+\Hh l3)
        +2\Hh{k+l}3(\Hh k2+\Hh l2)+2\Hh{k+l}2(\Hh k3+\Hh l3)
        -2\Hh k5-2\Hh l5\\
       &\quad-6S_{1,4}(k)-6S_{1,4}(l)-2S_{2,3}(k)-2S_{2,3}(l)
        +6U_{1,4}(k,l)+6U_{1,4}(l,k)+2U_{2,3}(k,l)+2U_{2,3}(l,k),
\end{align*}
with $U_{r,m}(a,b)$ abbreviated $U_{r,m}$ evaluated at $(a,b)$.
\end{proposition}

\begin{proof}[Proof sketch, with the full reduction recorded]
On the triangle $0<u<v<1$ substitute $u=xv$; both denominators expand
geometrically, and every term is a product of two beta-type integrals in $x$
and $v$ carrying powers of $\log$.  The termwise integration is justified by
monotone convergence after splitting the logarithms into signs, or by
integrating over $[\delta,1-\delta]^2$ and letting $\delta\to0$; both
triangle contributions have integrable boundary limits.  The result is a
bilinear combination of the tails
\[
 Z_m(N)=\sum_{t\ge N}t^{-m}\ (m\ge2),\qquad
 Z_1\ \text{regularised},\qquad
 \sum_{t\ge1}\frac1{t^{r}(t+h)^{m}},\qquad
 \sum_{t\ge1}\frac{\Hh{t+d-1}{m}}{t^{r}},
\]
with $h\le k_2$, $d=k_2+1$ and offsets governed by $k_1$ (and the mirror
triangle swaps $(k_1,s_1)\leftrightarrow(k_2,s_2)$).  The two divergent
$Z_1$-constants cancel pairwise ($a_1+b_1=0$ in each partial-fraction
splitting), so no regularisation ambiguity remains: every step is a limit of
finite rational identities.  Shifted products
$\sum_t t^{-r}(t+h)^{-m}$ reduce by partial fractions to zeta values and
harmonic numbers at $h$; the coupled tails reduce to the double zeta values
$\zeta(m,r)=\sum_{t>u\ge1}t^{-m}u^{-r}$ plus the finite sums $S,U$; and the
weight-at-most-five double zetas reduce by
\[
 \zeta(2,1)=\zeta(3),\quad \zeta(3,1)=\tfrac{\zeta(4)}4,\quad
 \zeta(2,2)=\tfrac{3\zeta(4)}4,\quad
 \zeta(4,1)=2\zeta(5)-\zeta(2)\zeta(3),
\]
\[
 \zeta(3,2)=3\zeta(2)\zeta(3)-\tfrac{11}2\zeta(5),\quad
 \zeta(2,3)=\tfrac92\zeta(5)-2\zeta(2)\zeta(3),
\]
together with $\sum_t H_t/t^2=2\zeta(3)$, $\sum_t H_t/t^3=\tfrac54\zeta(4)$
and $\zeta(2)^2=\tfrac52\zeta(4)$.  Since the $\zeta(5)$- and
$\zeta(2)\zeta(3)$-carrying heads $\zeta(4,1),\zeta(3,2),\zeta(2,3)$ enter
with $(k_1,k_2)$-independent coefficients, the top-weight coefficients are
constants, fixed by the calibration $I^{(2,2)}_{0,0}=I_0$.  Carrying out the
finite bookkeeping for each $(s_1,s_2)$ yields the displayed tables; the
tables were additionally re-derived by an independent implementation of the
same reduction and cross-checked in exact arithmetic on $0\le k,l<14$
(both the derived and the displayed compact forms), so each entry is an
identity of finite expressions whose verification is mechanical.

Two representative entries, worked in full, exhibit the bookkeeping.
Write $a=k_1$, $b=k_2$, $A=a+1$, $d=b+1$; set
$Z_m(N)=\sum_{t\ge N}t^{-m}$ for $m\ge2$ and $Z_1(N)=-H_{N-1}$ (the
regularised value; the divergent constants cancel pairwise),
$U_{r,m}(h)=\sum_{t\ge1}t^{-r}(t+h)^{-m}$,
$S_{r,m}(A,d)=\sum_{t\ge A}t^{-r}Z_m(t+d)$, and
$E_2=\sum_tH_t/t^2=2\zeta(3)$, $E_3=\sum_tH_t/t^3=\tfrac54\zeta(4)$
(these enter through $S_{r,1}$).  Let $F^{s_1,s_2}_{a,b}$ denote
the contribution of the triangle $0<u<v<1$, so that
$I^{(s_1,s_2)}_{k_1,k_2}
=(-1)^{s_1+s_2}\bigl[F^{s_1,s_2}_{a,b}+F^{s_2,s_1}_{b,a}\bigr]
/\bigl(\Gamma(s_1)\Gamma(s_2)\bigr)$
and
\[
 F^{s_1,s_2}_{a,b}
 =(-1)^{s_1+s_2-1}\sum_{i=0}^{s_1}\binom{s_1}i\,i!\,(s_1{-}i{+}s_2{-}1)!\,
   S_{i+1,\,s_1-i+s_2}(A,d)
 +(-1)^{s_1+s_2}\sum_{i=0}^{s_1-1}\binom{s_1{-}1}i(i{+}1)!\,(s_1{+}s_2{-}2{-}i)!\,
   Z_{i+2}(A)\,Z_{s_1+s_2-1-i}(A{+}d{-}1).
\]

\emph{$[\zeta(4)]$ of $I^{(1,2)}$.}
For $(s_1,s_2)=(1,2)$:
$F^{1,2}=2S_{1,3}+S_{2,2}-Z_2(A)Z_2(A{+}d{-}1)$, whose $\zeta(4)$ content is
$2\cdot\tfrac14+\tfrac34-\tfrac52=-\tfrac54$
(from $\zeta(3,1)$, $\zeta(2,2)$ and $\zeta(2)^2$; the finite sums $U_{r,m}(h)$
carry no $\zeta(4)$ since their zeta coefficients all come with positive
powers of $h^{-1}$).  For the mirror $(s_1,s_2)=(2,1)$:
$F^{2,1}=2S_{1,3}+2S_{2,2}+2S_{3,1}-Z_2Z_2-2Z_3Z_1$, and
$S_{3,1}$ contributes $-E_3=-\tfrac54\zeta(4)$; the $\zeta(4)$ content is
$\tfrac12+\tfrac32-\tfrac52-\tfrac52=-3$.  Hence
$[\zeta(4)]I^{(1,2)}=(-1)^{1+2-2}\bigl(-\tfrac54-3\bigr)/\bigl(\Gamma(1)\Gamma(2)\bigr)
=\tfrac{17}4$, as tabulated --- note that without the
Laplace sign the value would come out as $-\tfrac{17}4$.

\emph{$[\zeta(2)]$ of $I^{(2,2)}$.}
For $(s_1,s_2)=(2,2)$:
$F^{2,2}=-\bigl[6S_{1,4}+4S_{2,3}+2S_{3,2}\bigr]
+2Z_2(A)Z_3(A{+}d{-}1)+2Z_3(A)Z_2(A{+}d{-}1)$.
The isolated-$\zeta(2)$ content of the $U$-sums is
$[\zeta(2)]U_{1,4}(h)=-h^{-3}$, $[\zeta(2)]U_{2,3}(h)=3h^{-3}$,
$[\zeta(2)]U_{3,2}(h)=-3h^{-3}$, so summing over $h\le b$,
$[\zeta(2)]S_{1,4}=\Hh b3$, $[\zeta(2)]S_{2,3}=-3\Hh b3$,
$[\zeta(2)]S_{3,2}=3\Hh b3-\Hh a3$ (the last term from
$-\sum_{t\le a}t^{-3}\zeta(2)$ inside the coupled tail), while the
$Z$-products contribute $-2\Hh{a+b}3-2\Hh a3$.  Altogether
$[\zeta(2)]F^{2,2}_{a,b}=-\bigl[6-12+6\bigr]\Hh b3+2\Hh a3-2\Hh{a+b}3-2\Hh a3
=-2\Hh{a+b}3$, and adding the mirror gives
$[\zeta(2)]I^{(2,2)}=-4\Hh{k+l}3$.
\end{proof}

\begin{remark}
The finite Euler sums are not an artefact: the pair
$(S_{r,m},U_{r,m})$ is exactly the local trace the sine kernel leaves behind
after the two-triangle integration, and the \emph{basis} in which one writes
them matters.  The expressions above use the fixed basis
$\{$products of $\Hh\cdot\cdot$ at $k,l,k+l$;\ $S_{r,m}$ at $k,l$;\
$U_{r,m}$ at $(k,l),(l,k)\}$; different-looking outputs of the same
reduction are equal modulo the finite quasi-shuffle relations, e.g.\
$S_{r,m}(x)+S_{m,r}(x)=\Hh xr\Hh xm+\Hh x{r+m}$ and
$U_{r,m}(a,b)=\Hh ar\Hh{a+b}m-\sum_{j=1}^{a}\Hh{j-1}r(b+j)^{-m}$.
\end{remark}

\section{The four kernel identities}\label{sec:kernelids}

Throughout this section we use the abbreviations
$A_{kl}=C_{22}$, $B_{kl}=C_{12}$, $C_{kl}=C_{21}$, $D_{kl}=C_{11}$ of
Proposition~\ref{prop:pf} (so $B_{kl}=A_{kl}L_k$ etc.), and the one-variable
sections of $R:=R_n(x-n-1,\,y-n-1)$, i.e.\ of
\[
 R(x,y)=\frac{P(x)P(y)P(x+y)}{Q(x)^2Q(y)^2},\qquad
 P(z)=\prod_{r=1}^n(z-r),\quad Q(z)=\prod_{r=0}^n(z+r),
\]
whose partial fractions on the grid $x=-k$, $y=-l$ carry the same data
$A,B,C,D$.  For fixed $l$ let
\[
 g_l(x)=\lim_{y\to-l}(y+l)^2R(x,y)
 =\sum_{k=0}^n\Bigl[\frac{A_{kl}}{(x+k)^2}+\frac{B_{kl}}{x+k}\Bigr],
\qquad
 q_l(x)=\frac{\partial}{\partial y}\Bigl[(y+l)^2R(x,y)\Bigr]_{y=-l}
 =\sum_{k=0}^n\Bigl[\frac{C_{kl}}{(x+k)^2}+\frac{D_{kl}}{x+k}\Bigr].
\]
From the product form,
\begin{equation}\label{eq:gq}
 g_l(x)=c_l\,\frac{P(x)P(x-l)}{Q(x)^2},
 \qquad
 q_l(x)=\frac{P(x)}{Q(x)^2}\bigl[c'_l\,P(x-l)+c_l\lambda_l P(x-l)\bigr]',
\end{equation}
more precisely $q_l(x)=P(x)Q(x)^{-2}\bigl(\mu_l P'(x-l)+\nu_l P(x-l)\bigr)$
for finite constants $\mu_l,\nu_l$ (the $y$-derivative never touches the
factor $P(x)$).  Consequently:
\begin{itemize}
\item[(Z1)] $g_l(j)=0$ for $1\le j\le n+l$ \ (both numerator products
vanish on $1\le j\le n$ resp.\ $l+1\le j\le n+l$; the range is sharp);
\item[(Z2)] $g_l$ has a \emph{double} zero on $l<j\le n$, so also
$g_l'(j)=0$ there (sharp);
\item[(Z3)] $q_l(j)=0$ for $1\le j\le n$ \ (the surviving factor $P(j)=0$;
sharp, and the zero has order exactly one).
\end{itemize}
Moreover $g_l,q_l=O(x^{-2})$ as $x\to\infty$, so the sums of their residues
vanish:
\begin{equation}\label{eq:res0}
 \sum_{k=0}^nB_{kl}=0,
 \qquad
 \sum_{k=0}^nD_{kl}=0
 \qquad\text{for each fixed } l,
\end{equation}
and symmetrically with the roles of $k$ and $l$ exchanged
($A_{kl}=A_{lk}$, $C_{kl}=B_{lk}$, $D_{kl}=D_{lk}$).

\begin{proposition}[$\zeta(5)$, $\zeta(2)\zeta(3)$, $\zeta(4)$]
\label{prop:z5z4}
For every $n$:
\[
 [\zeta(5)]\sum_{k,l}T\,W_B=2Q_n,\qquad
 [\zeta(2)\zeta(3)]\sum_{k,l}T\,W_B=4Q_n,\qquad
 [\zeta(4)]\sum_{k,l}T\,W_B=0 .
\]
\end{proposition}

\begin{proof}
The first two are immediate from Proposition~\ref{prop:ratparts} and
\eqref{eq:Qn}: only $I^{(2,2)}$ carries these values, with constant
coefficients $2$ and $4$.

For $\zeta(4)$: by Proposition~\ref{prop:ratparts} the coefficient is
$\frac{17}4\sum_{k,l}T\,(L_k+L_l)$.  For fixed $l$,
$\sum_kT\,L_k$ is proportional to $\sum_kB_{kl}=0$ by \eqref{eq:res0};
summing over $l$, and treating $L_l$ symmetrically, gives zero.
\end{proof}

\begin{proposition}[$\zeta(3)$]\label{prop:z3}
For every $n$, \ $[\zeta(3)]\sum_{k,l}T\,W_B=0$.
\end{proposition}

\begin{proof}
By Proposition~\ref{prop:ratparts}, up to the common factor
$(-1)^nn!^{-1}$,
\[
 [\zeta(3)]\sum T\,W_B
 \;\propto\;
 2\sum_{k,l}\Bigl[
 A_{kl}\bigl(\Hh k2{+}\Hh l2{-}2\Hh{k+l}2\bigr)
 +B_{kl}(H_k{-}H_{k+l})
 +C_{kl}(H_l{-}H_{k+l})
 +D_{kl}\Bigr].
\]
The $D$-term vanishes by \eqref{eq:res0}.  Using
$H_k-H_{k+l}=-\sum_{j=1}^{l}\frac1{k+j}$,
$\Hh k2-\Hh{k+l}2=-\sum_{j=1}^{l}\frac1{(k+j)^2}$ and the
$k\leftrightarrow l$ symmetry to fold the $\Hh l2$- and $C$-halves onto the
$k$-side, the bracket becomes
\[
 -2\sum_{l=0}^n\sum_{j=1}^{l}\sum_{k=0}^n
 \Bigl[\frac{A_{kl}}{(k+j)^2}+\frac{B_{kl}}{k+j}\Bigr]
 =-2\sum_{l=0}^n\sum_{j=1}^{l}g_l(j)=0
\]
by (Z1), since $1\le j\le l\le n+l$.
\end{proof}

\begin{proposition}[$\zeta(2)$]\label{prop:z2}
For every $n$,
\[
 [\zeta(2)]\sum_{k,l}T\,W_B=-4\sum_{k,l}T\,w_3^{\mathrm{sym}},
 \qquad
 w_3^{\mathrm{sym}}
 =\tfrac12\bigl(\Hh{n+k}3+\Hh{n+l}3\bigr)
 -\tfrac\Psi2\bigl(\Hh{n+k}2-\Hh{n+l}2\bigr).
\]
\end{proposition}

\begin{proof}
Step 1.  By Proposition~\ref{prop:ratparts} (note
$[\zeta(2)]I^{(2,2)}=-4\Hh{k+l}3$), and again up to the common nonzero
factor,
\[
 [\zeta(2)]\sum T\,W_B
 \propto
 \sum_{k,l}\Bigl[-4A_{kl}\Hh{k+l}3
 +B_{kl}\bigl(\Hh l2{-}2\Hh{k+l}2\bigr)
 +C_{kl}\bigl(\Hh k2{-}2\Hh{k+l}2\bigr)
 +D_{kl}\bigl(H_k{+}H_l{-}2H_{k+l}\bigr)\Bigr].
\]
We first claim this equals
$\sum_{k,l}\bigl[-4A_{kl}\Hh{k+l}3-(B_{kl}+C_{kl})\Hh{k+l}2\bigr]$.
The difference is
\[
 \sum_{k,l}\Bigl[B_{kl}\bigl(\Hh l2{-}\Hh{k+l}2\bigr)
 +C_{kl}\bigl(\Hh k2{-}\Hh{k+l}2\bigr)
 +D_{kl}\bigl(H_k{-}H_{k+l}\bigr)+D_{kl}\bigl(H_l{-}H_{k+l}\bigr)\Bigr],
\]
and unfolding the harmonic differences into ranges as in
Proposition~\ref{prop:z3}, then folding the $(C,D_{(l\text{-side})})$-half
onto the $k$-side by symmetry, the difference equals
\[
 -2\sum_{k=0}^n\sum_{j=1}^{k}\sum_{l=0}^n
 \Bigl[\frac{B_{kl}}{(j+l)^2}+\frac{D_{kl}}{j+l}\Bigr]
 =-2\sum_{k}\sum_{j=1}^{k}q^{\,\mathrm{mirror}}_k(j)=0
\]
by the mirror of (Z3) (the simple-pole-in-$x$ section, as a function of
$y$, vanishes at $y=j$ for $1\le j\le n$, and $j\le k\le n$).

Step 2.  It remains to identify
$\sum T\bigl[\Hh{k+l}3+\tfrac14(L_k+L_l)\Hh{k+l}2\bigr]$
with $\sum T\,w_3^{\mathrm{sym}}$; multiplying by $-4$ then gives the claim.
Multiply $w_3^{\mathrm{sym}}$ by $4A_{kl}$ and use $L_k-L_l=-2\Psi$ and
symmetry: the difference of the two summands reduces to
\[
 \sum_{k,l}\Bigl[4A_{kl}\bigl(\Hh{n+k}3-\Hh{k+l}3\bigr)
 +2B_{kl}\bigl(\Hh{n+k}2-\Hh{n+l}2-\Hh{k+l}2\bigr)\Bigr]
\]
after folding by symmetry.  For fixed $l$ the term $-2B_{kl}\Hh{n+l}2$ sums
to zero by \eqref{eq:res0}.  Expanding the two remaining harmonic
differences over the range $k+l<i\le n+k$, i.e.\ $i=k+j$ with $l<j\le n$,
leaves
\[
 \sum_{l=0}^n\sum_{j=l+1}^{n}\sum_{k=0}^n
 \Bigl[\frac{4A_{kl}}{(k+j)^3}+\frac{2B_{kl}}{(k+j)^2}\Bigr]
 =-2\sum_{l=0}^n\sum_{j=l+1}^{n}g_l'(j)=0
\]
by (Z2).
\end{proof}

\section{The middle row}\label{sec:middle}

This section proves Theorem~\ref{thm:middle}.  It is logically independent
of Sections~\ref{sec:recomb}--\ref{sec:kernelids} except for reusing
(Z1)--(Z2) and \eqref{eq:res0}; that independence is what powers the
subtraction argument of Section~\ref{sec:anchor}.

Brown--Zudilin's descent \cite[eq.~(I3)]{BZ}, specialised to the totally
symmetric parameters (write $k=n+l$ in their summation index), reads
\begin{equation}\label{eq:I3}
 2I''_n=\sum_{l=0}^{n}(-1)^l\binom{n+l}{n}\binom nl^2
 \,J_3\bigl(n,n,n,n{-}l;\,n,n,n{+}l\bigr),
\end{equation}
where $J_3$ is the generalised Beukers integral and
$I''_n=Q_n\zeta(3)-\hat P_n$.  (The factor $2$ is discussed below; it is
forced uniformly in $n$, not fitted.)  Each term satisfies the Rhin--Viola
condition $(n-l)+(n+l)=n+n$, so by \cite[Sec.~3]{Zu02} it is a linear form
$2A\zeta(3)-B$ whose coefficients come from the partial fractions of an
explicit rational function: using Zudilin's parameter correspondence
(\cite[eqs.~(3.8),(3.10)]{Zu02}), the relevant function for the term $l$ is
\[
 R_{n,l}(t)=
 \frac{\prod_{i=1}^n(t+i)\,\prod_{i=1}^n(t+i-l)}
      {\prod_{i=n+1}^{2n+1}(t+i)^2},
\]
with double poles at $t=-n-1-k$, $0\le k\le n$.  Writing
\[
 R_{n,l}(t)=\sum_{k=0}^n\Bigl[\frac{A_{kl}}{(t+n+1+k)^2}
 +\frac{B_{kl}}{t+n+1+k}\Bigr],
\]
the same product evaluations as in Proposition~\ref{prop:pf} give
\begin{equation}\label{eq:AB}
 A_{kl}=\binom{n+k}{n}\binom{n+k+l}{n}\binom nk^2,
 \qquad
 B_{kl}=A_{kl}\,L_k .
\end{equation}
Zudilin's Lemma~4 \cite[eq.~(2.13) and its proof]{Zu02} evaluates the
corresponding series as
\[
 J_3(n,n,n,n{-}l;n,n,n{+}l)
 =(-1)^l\Bigl\{2\sum_kA_{kl}\,\zeta(3)
 -\sum_kA_{kl}\bigl[2\Hh{k+l}3+L_k\Hh{k+l}2\bigr]\Bigr\}.
\]
Substituting into \eqref{eq:I3} and using that the outer coefficient times
$A_{kl}$ is exactly $T(n,k,l)$ (the six binomials merely reorder), we obtain
the real-number identity
\begin{equation}\label{eq:middlereal}
 2Q_n\zeta(3)-2\hat P_n
 =2\sum_{k,l}T\,\zeta(3)
 -\sum_{k,l}T\bigl[2\Hh{k+l}3+L_k\Hh{k+l}2\bigr].
\end{equation}
The $\zeta(3)$ terms on the two sides are \emph{identically equal
rationals} multiplied by the same real number $\zeta(3)$, by \eqref{eq:Qn};
subtracting them (no independence needed) leaves
\begin{equation}\label{eq:M}
 \hat P_n=\sum_{k,l}T\Bigl[\Hh{k+l}3+\tfrac12L_k\Hh{k+l}2\Bigr]
 =\sum_{k,l}T\Bigl[\Hh{k+l}3+\tfrac14(L_k{+}L_l)\Hh{k+l}2\Bigr],
\end{equation}
the second equality by $k\leftrightarrow l$ symmetry of $T$.  Step 2 of
Proposition~\ref{prop:z2} identifies the right side with
$\sum T\,w_3^{\mathrm{sym}}$.  Finally $w_3^{\mathrm{sym}}$ is the
symmetrisation of $\hat w_3$: their difference is $k\leftrightarrow l$
antisymmetric, and $T$ is symmetric, so
$\sum T\,w_3^{\mathrm{sym}}=\sum T\,\hat w_3$.  This proves
Theorem~\ref{thm:middle}, up to two normalisation points:

\emph{The factor $2$ in \eqref{eq:I3}.}  The iterated-residue derivation of
(I3) in \cite{BZ} produces the left-hand side as the residue integral
displayed there; comparing its $\zeta(3)$ coefficient, computed by the
right-hand side of \eqref{eq:middlereal} to be $2\sum T=2Q_n$, with the
$\zeta(3)$ coefficient $Q_n$ of $I''_n$ shows the displayed residue object
equals $2I''_n$ for every $n$.  Both coefficients are explicit rationals
inside one derivation, so this comparison involves no independence
statement.  (At $n=0$ it is the classical calibration: the residue integral
is Beukers' $2\zeta(3)$ while $I''_0=\zeta(3)$.)

\emph{$I''_n$ versus \eqref{eq:In}.}  The identity
$I''_n=Q_n\zeta(3)-\hat P_n$ is \cite{BZ}'s definition of the middle-row
pair via the same recursion-plus-initial-values argument as \eqref{eq:In};
its role here is only to name the rational number $\hat P_n$.
\hfill$\qed$

\section{The subtraction anchor: proof of Theorem~\ref{thm:anchor}}
\label{sec:anchor}

By Propositions~\ref{prop:recomb} and~\ref{prop:ratparts},
\begin{equation}\label{eq:barnesside}
 I_n=\sum_{k,l}T\,W_B
 =2Q_n\zeta(5)+4Q_n\zeta(2)\zeta(3)+0\cdot\zeta(4)+0\cdot\zeta(3)
 -4\Bigl[\sum_{k,l}T\,w_3^{\mathrm{sym}}\Bigr]\zeta(2)
 +\sum_{k,l}T\,W^{(1)},
\end{equation}
using Propositions~\ref{prop:z5z4}--\ref{prop:z2} for the five upper
coefficients; here $W^{(1)}=[1]W_B$ is \eqref{eq:W1}.  On the other hand
\eqref{eq:In} gives
\[
 I_n=2Q_n\zeta(5)+4Q_n\zeta(2)\zeta(3)-4\hat P_n\zeta(2)-2P_n .
\]
Subtract: by Theorem~\ref{thm:middle},
$\sum T\,w_3^{\mathrm{sym}}=\hat P_n$, so every transcendental term cancels
\emph{exactly}, leaving the equality of real numbers
\[
 \sum_{k,l}T\,W^{(1)}=-2P_n .
\]
This is Theorem~\ref{thm:anchor}.\hfill$\qed$

\begin{remark}
Theorem~\ref{thm:anchor} is an unconditional closed form for the top row:
$P_n$ equals an explicit double sum of products of binomials, harmonic
numbers, and the finite Euler sums $S_{r,m},U_{r,m}$.  It is the exact
weight-five analogue of the middle row's \eqref{eq:M}, and it is the first
representation of $P_n$ obtained without any coefficient comparison.
\end{remark}

\begin{remark}
The chain uses \eqref{eq:In} and \eqref{eq:Qn} from \cite{BZ}, whose proofs
rest on a certified creative-telescoping recursion for $I_n$ plus exact
evaluation of $I_0,I_1,I_2$; an independent certificate for the same
recursion, in normalised form, is part of the present programme's
repository.  No other external input enters.
\end{remark}

\section{The remaining bridge to the compact weight}\label{sec:T3}

Theorems~\ref{thm:anchor} and~\ref{thm:top} differ by the single identity
\eqref{eq:T3}, equivalently
\[
 \Delta_n:=\sum_{k,l=0}^nT(n,k,l)\,
 \bigl[W^{(1)}(k,l)+2w_5^{\mathrm{sym}}(n,k,l)\bigr]=0
 \qquad(n\ge0).
\]
This is a statement about finite sums only; no integrals or zeta values
remain.  We record its status precisely.

\subsection{What is verified}
$\Delta_n=0$ has been checked in exact rational arithmetic for $n\le12$
(all cells, all six basis coefficients of the underlying decomposition
compared independently), and modulo two 31-bit primes for $n\le80$.
Moreover $P_n=\sum T\,w_5$ itself has been verified against the exact
recursion ladder for $n\le360$ at two primes.

\subsection{The proof strategy, and its measured state}
All identities of Section~\ref{sec:kernelids} were proved by exhibiting the
relevant weighted sum as a combination of \emph{null functionals}: sums
that vanish for every $n$ because of (Z1)--(Z3), \eqref{eq:res0}, the
two-variable vanishing $R(x,y)=0$ at lattice points and along the moving
anti-diagonals $x+y=m$ ($1\le m\le n$, all Laurent coefficients), and the
one-variable residue lemma applied to $R_k(z)\rho(z)$ for admissible
rational $\rho$ (``pole-raising jets'', which carry the higher logarithmic
jets $e_2,\dots,e_5$ of the local germ that fixed-pole evaluations cannot
see).  The natural certificate for \eqref{eq:T3} would be the same shape,
in its \emph{folded} form: since $B_{kl}=A_{kl}L_k$ and
$D_{kl}=A_{kl}(L_kL_l-C_2)$ with $L$'s polynomial in harmonic numbers, any
cellwise identity
\[
 \bigl[W^{(1)}+2w_5^{\mathrm{sym}}\bigr](n,k,l)
 =\sum_i x_i\,u_i(n,k,l),
 \qquad u_i\ \text{null},
\]
implies \eqref{eq:T3}.

We record the measured state of this reduction (all statements modulo a
31-bit prime, with the protocol: solve on $n\le N$, test consistency on the
dependent rows, hold out $n=N{+}1$).  With the generator inventory
consisting of: weighted fixed-pole evaluations (including inverse-power and
$n$-coupled weights), derivative evaluations on the double-zero range,
anti-diagonal and vertical Laurent band families, pole-raising jets to
order five, weighted value towers
$\sum_{j\in\mathrm{range}}\varphi(j)\,(R_k\rho)(-j)=0$ with $k$-side letter
multipliers (these produce the coupled sums $\sum_j\varphi(j)/(j+l)^s$ up
to $s=5$ that the $U$-letters of $W^{(1)}$ require), and two-variable
iterated-residue jets $\sum_{k,l}\Res_y\Res_x[R\,\rho(x)\sigma(y)]=0$ with
lattice pole-raising in both variables (these turn out, measurably, to lie
in the per-fibre span, so they enlarge nothing) ---
about $5{,}600$ null functionals in all, each individually proved and
machine-calibrated --- the three-component system is consistent through
$n\le14$ (with over a thousand dependent rows) but \emph{inconsistent} at
$n\le17$, and the folded system is inconsistent at $n\le20$ (514 of 3310
rows).  A further family of \emph{inhomogeneous moment towers} --- keeping
the residue at infinity,
$\sum_{l}\Res_{z=l}[z^mR_k\rho]=[z^{-m-1}](R_k\rho)\big|_\infty$, whose
densities carry polynomial $l$-weights and lower jet coefficients, with the
explicit right-hand sides constrained by their own cost rows --- adds
$726$ further independent directions and closes the folded system through
$n\le20$, but the properly determined test (all cells $n\le25$, more rows
than columns) is again inconsistent, with the rank of the entire calculus
saturating near $3500$ while the deficit grows linearly in the number of
rows.

The conclusion is a structural negative rather than a missing family:
\emph{no cellwise certificate with fixed letter coefficients over the
local/residue calculus expresses \eqref{eq:T3}}.  This matches a formal
obstruction: the folded target contains $l$-free nested monomials
($S_{1,4}(k)$, $S_{2,3}(k)$, \dots) which no per-fibre or iterated-residue
functional can produce monomially.  The certificate for \eqref{eq:T3} must
therefore live where every proved identity of this type lives: in the
creative-telescoping class, with coefficients \emph{rational} in $(n,k,l)$,
\[
 T\cdot\bigl[W^{(1)}+2w_5^{\mathrm{sym}}\bigr]
 =\Delta_k R+\Delta_l S,
\]
or its two-stage one-variable form (an $l$-telescoper for the fibre sums
$V(n,k)=\sum_lT\cdot[\,\cdot\,]$, then a $k$-telescoper for $\sum_kV=0$).
This is well-founded: the $U$-letters have rational $l$-shift closure over
the bare alphabet
($U_{r,m}(k,l{+}1)-U_{r,m}(k,l)=\sum_{t\le k}t^{-r}(t{+}l{+}1)^{-m}$
partial-fractions into rationals times bare harmonic differences), so the
defect's shift module is finite-rank over $\QQ(n,k,l)$ and the telescoper
is the kernel of a finite matrix.  That computation is bounded and is the
remaining step.

\subsection{Why the compact weight is expected}
The companion note \emph{The Laurent--jet origin of the compact
Brown--Zudilin weights} shows that $\hat w_3$ and $w_5$ are forced as
Bell-polynomial coefficients of the anti-diagonal logarithmic jets of the
local germ of $R_n$ at its finite double poles --- the direction that
annihilates the coupled factor $P(x+y)$.  Identity \eqref{eq:T3} is the
statement that the same jets exhaust the rational content of the four
universal integrals after $T$-weighted summation; every piece of
$W^{(1)}$'s nested content lies in the image of the weighted-tower
functionals, which is the structural reason the bridge is expected to close
over the proved generator space.

\begin{thebibliography}{9}
\bibitem{BZ} F.~Brown, W.~Zudilin,
\emph{On cellular rational approximations to $\zeta(5)$},
arXiv:2210.03391.
\bibitem{Zu02} W.~Zudilin,
\emph{Arithmetic of linear forms involving odd zeta values},
J.\ Th\'eor.\ Nombres Bordeaux 16 (2004), 251--291.
\bibitem{Be79} F.~Beukers, \emph{A note on the irrationality of $\zeta(2)$
and $\zeta(3)$}, Bull.\ London Math.\ Soc.\ 11 (1979), 268--272.
\bibitem{RV01} G.~Rhin, C.~Viola, \emph{The group structure for $\zeta(3)$},
Acta Arith.\ 97 (2001), 269--293.
\end{thebibliography}

\end{document}
